% =============================================================
%  AC power transmission — the general representation
%  Applies to overhead lines and cables alike; the type-specific
%  parameters are in 04-overhead-lines.tex and 05-cables.tex, the
%  worked numbers in 06-reference-line.tex.
%  Figures: nominal pi-section in the per-unit-length parameters and
%  as a two-port, both inline TikZ in the report colours.
%  Slides derived from this section:
%  presentation/design/ac-transmission/
% =============================================================
\section{AC Power Transmission}
\label{sec:ac-chapter}

The familiar statement that active power flows on the angle between two
voltages, $P = U_S \cdot U_R / X \cdot \sin\delta$, is a special case. It
holds for a short, lossless line, and the error it carries grows with length
until it is useless. This section builds the representation from the line
parameters upwards: the four per-unit-length quantities, the two-port that
follows from them, the exact power transfer relation, the conditions under
which it collapses to the simple form, and the distributed solution that
replaces the lumped one on long lines.

Nothing in this section distinguishes an overhead line from a cable. The two
differ only in the values of the four parameters, and in what those values
make critical: Section~\ref{sec:ohl} covers the overhead line,
Section~\ref{sec:cable} the cable, and Section~\ref{sec:reference-line}
carries the whole chain through on one worked example.

\subsection{Fundamentals}
\label{sec:ac-transmission}

The representation below is general: it applies to an overhead line
and to a cable alike, and the rest of this chapter specialises it.

\subsubsection{Line parameters}
\label{sec:ac-parameters}

A transmission line is described by four parameters per unit length, two in
series with the conductor and two in shunt to the return path. Each
corresponds to one physical effect, Table~\ref{tab:ac-parameters}.

\begin{table}[!ht]
  \centering
  \begin{tabularx}{\linewidth}{@{}l l X@{}}
    \toprule
    & \multicolumn{1}{c}{\textbf{Parameter}} & \textbf{Physical effect} \\
    \midrule
    $R'$ & Resistance  & Ohmic loss in the conductors \\
    $L'$ & Inductance  & Magnetic field around the conductors \\
    $G'$ & Conductance & Leakage through the insulation and corona \\
    $C'$ & Capacitance & Electric field to the other conductors and to earth \\
    \bottomrule
  \end{tabularx}
  \caption{Line parameters per unit length}
  \label{tab:ac-parameters}
\end{table}

They enter as a series impedance and a shunt admittance per unit length,
\begin{equation}
  \underline{Z}' = R' + j \cdot \omega \cdot L'
  \qquad
  \underline{Y}' = G' + j \cdot \omega \cdot C'
  \label{eq:ac-zy-prime}
\end{equation}

$R'$ and $L'$ follow from the conductor itself and from the arrangement of
the conductors; $G'$ and $C'$ from the insulation between them. The
expressions differ between an overhead line and a cable and are given in
Sections~\ref{sec:ohl} and~\ref{sec:cable}.

Taken over a finite length the four parameters arrange themselves as the
nominal $\pi$-section of Figure~\ref{fig:ac-pi-section}: the series branch
carries $R'$ and $L'$, and the shunt branch $G'$ and $C'$ is split equally
between the two ends.

\begin{figure}[!ht]
  \centering
  \begin{tikzpicture}[
      x=0.60cm, y=0.60cm, line width=0.8pt,
      every node/.style={inner sep=1.5pt, font=\small},
      rail/.style={atFg, line width=0.8pt},
      box/.style={draw=atFg, fill=atBgSubtle, line width=0.8pt,
                  minimum width=1.0cm, minimum height=0.46cm},
      vbox/.style={draw=atFg, fill=atBgSubtle, line width=0.8pt,
                   minimum width=0.60cm, minimum height=1.05cm},
      term/.style={circle, draw=atFg, fill=white, line width=0.8pt,
                   minimum size=3.8pt, inner sep=0pt},
      dim/.style={dotted, atFgMuted, line width=0.9pt},
      ext/.style={<->, >=Latex, atFgMuted, line width=0.8pt},
    ]
    % series branch
    \draw[rail] (0,2.8) -- (2.77,2.8);
    \node[box] at (3.6,2.8) {$R'$};
    \draw[rail] (4.43,2.8) -- (4.77,2.8);
    \node[box] at (5.6,2.8) {$L'$};
    \draw[rail] (6.43,2.8) -- (9.2,2.8);
    % return conductor
    \draw[rail] (0,0) -- (9.2,0);
    % left shunt leg
    \draw[rail] (1.3,2.8) -- (1.3,2.4);
    \draw[rail] (0.7,2.4) -- (1.9,2.4);
    \draw[rail] (0.7,2.4) -- (0.7,2.28);
    \draw[rail] (1.9,2.4) -- (1.9,2.28);
    \node[vbox] at (0.7,1.4) {$\frac{G'}{2}$};
    \node[vbox] at (1.9,1.4) {$\frac{C'}{2}$};
    \draw[rail] (0.7,0.52) -- (0.7,0.4);
    \draw[rail] (1.9,0.52) -- (1.9,0.4);
    \draw[rail] (0.7,0.4) -- (1.9,0.4);
    \draw[rail] (1.3,0.4) -- (1.3,0);
    % right shunt leg
    \draw[rail] (7.9,2.8) -- (7.9,2.4);
    \draw[rail] (7.3,2.4) -- (8.5,2.4);
    \draw[rail] (7.3,2.4) -- (7.3,2.28);
    \draw[rail] (8.5,2.4) -- (8.5,2.28);
    \node[vbox] at (7.3,1.4) {$\frac{G'}{2}$};
    \node[vbox] at (8.5,1.4) {$\frac{C'}{2}$};
    \draw[rail] (7.3,0.52) -- (7.3,0.4);
    \draw[rail] (8.5,0.52) -- (8.5,0.4);
    \draw[rail] (7.3,0.4) -- (8.5,0.4);
    \draw[rail] (7.9,0.4) -- (7.9,0);
    % terminals
    \node[term] at (0,2.8) {};    \node[term] at (0,0) {};
    \node[term] at (9.2,2.8) {};  \node[term] at (9.2,0) {};
    % length
    \draw[dim] (0,-0.3) -- (0,-1.6);
    \draw[dim] (9.2,-0.3) -- (9.2,-1.6);
    \draw[ext] (0,-1.2) -- (9.2,-1.2);
    \node[below=1pt, atFgMuted] at (4.6,-1.2) {$l$};
  \end{tikzpicture}
  \caption{Nominal $\pi$-section of a line of length $l$, written in the
           per-unit-length parameters}
  \label{fig:ac-pi-section}
\end{figure}

\subsubsection{The line as a two-port}
\label{sec:ac-two-port}

Over a length $l$ the parameters give a series impedance and a shunt
admittance
\begin{equation}
  \underline{Z} = \underline{Z}' \cdot l
  \qquad
  \underline{Y} = \underline{Y}' \cdot l
  \label{eq:ac-zy}
\end{equation}

\begin{figure}[!ht]
  \centering
  \begin{tikzpicture}[
      x=0.75cm, y=0.75cm, line width=0.8pt,
      every node/.style={inner sep=1.5pt, font=\small},
      rail/.style={atFg, line width=0.8pt},
      box/.style={draw=atFg, fill=atBgSubtle, line width=0.8pt,
                  minimum width=1.1cm, minimum height=0.46cm},
      vbox/.style={draw=atFg, fill=atBgSubtle, line width=0.8pt,
                   minimum width=0.55cm, minimum height=1.1cm},
      term/.style={circle, draw=atFg, fill=white, line width=0.8pt,
                   minimum size=3.8pt, inner sep=0pt},
      flow/.style={->, >=Latex, atFg, line width=0.8pt},
    ]
    % series branch
    \draw[rail] (0,2.6) -- (2.27,2.6);
    \node[box] at (3,2.6) {$\underline{Z}$};
    \draw[rail] (3.73,2.6) -- (6,2.6);
    % return conductor
    \draw[rail] (0,0) -- (6,0);
    % shunt legs
    \draw[rail] (1.2,2.6) -- (1.2,2.035);
    \node[vbox] at (1.2,1.3) {$\frac{\underline{Y}}{2}$};
    \draw[rail] (1.2,0.565) -- (1.2,0);
    \draw[rail] (4.8,2.6) -- (4.8,2.035);
    \node[vbox] at (4.8,1.3) {$\frac{\underline{Y}}{2}$};
    \draw[rail] (4.8,0.565) -- (4.8,0);
    % terminals
    \node[term] at (0,2.6) {};  \node[term] at (0,0) {};
    \node[term] at (6,2.6) {};  \node[term] at (6,0) {};
    % terminal currents
    \draw[flow] (-1.0,2.6) -- (-0.4,2.6);
    \node[left] at (-1.15,2.6) {$\underline{I}_S$};
    \draw[flow] (6.4,2.6) -- (7.0,2.6);
    \node[right] at (7.15,2.6) {$\underline{I}_R$};
    % terminal voltages, inset so they clear the terminals
    \draw[flow] (0,2.25) -- (0,0.35);
    \node[left=2pt] at (0,1.3) {$\underline{U}_S$};
    \draw[flow] (6,2.25) -- (6,0.35);
    \node[right=2pt] at (6,1.3) {$\underline{U}_R$};
  \end{tikzpicture}
  \caption{The same $\pi$-section as a two-port: the whole series impedance
           $\underline{Z}$ between the rails, half the shunt admittance
           $\underline{Y}$ at each end}
  \label{fig:ac-pi-zy}
\end{figure}

which the nominal $\pi$ arranges as in Figure~\ref{fig:ac-pi-zy}: the whole
series impedance between the rails and half the shunt admittance at each
end. This is the standard two-port description of a linear network, of which
the line is one case; Appendix~\ref{sec:two-port} gives the general form for
the reader who wants it, and is not needed here. The ports are named by
function throughout this section, the sending end and the receiving end. The
line is then a two-port,
\begin{equation}
  \begin{bmatrix} \underline{U}_S \\ \underline{I}_S \end{bmatrix}
  =
  \begin{bmatrix} \underline{A} & \underline{B} \\
                  \underline{C} & \underline{D} \end{bmatrix}
  \cdot
  \begin{bmatrix} \underline{U}_R \\ \underline{I}_R \end{bmatrix}
  \label{eq:ac-abcd}
\end{equation}
with, for the nominal $\pi$,
\begin{equation}
  \underline{A} = \underline{D} = 1 + \frac{\underline{Z} \cdot \underline{Y}}{2},
  \qquad
  \underline{B} = \underline{Z},
  \qquad
  \underline{C} = \underline{Y} \cdot
      \left(1 + \frac{\underline{Z} \cdot \underline{Y}}{4}\right)
  \label{eq:ac-abcd-pi}
\end{equation}

The shunt admittance shows itself in $\underline{A}$. A line that did not
charge itself would have $\underline{A} = 1$ and would behave as a single
series impedance; how far $\left|\underline{A}\right|$ departs from unity is
a measure of how far it is from that.

\subsubsection{Power transfer}
\label{sec:ac-power-transfer}

Solving~\eqref{eq:ac-abcd} for the receiving-end current gives
\begin{equation}
  \underline{I}_R
    = \frac{\underline{U}_S - \underline{A} \cdot \underline{U}_R}
           {\underline{B}}
  \label{eq:ac-ir}
\end{equation}
Taking the receiving-end voltage as the phase reference, so that $\delta$ is
the angle by which the sending-end voltage leads it,
\begin{equation}
  \underline{U}_S = U_S \cdot e^{\,j \cdot \delta},
  \quad
  \underline{U}_R = U_R,
  \quad
  \underline{A} = \left|A\right| \cdot e^{\,j \cdot \vartheta_A},
  \quad
  \underline{B} = \left|B\right| \cdot e^{\,j \cdot \vartheta_B}
  \label{eq:ac-reference}
\end{equation}
the receiving-end complex power $\underline{S}_R = \underline{U}_R \cdot
\underline{I}_R^{*}$ separates into
\begin{align}
  P_R &= \frac{U_S \cdot U_R}{\left|B\right|}
         \cdot \cos\left(\vartheta_B - \delta\right)
       - \frac{\left|A\right| \cdot U_R^2}{\left|B\right|}
         \cdot \cos\left(\vartheta_B - \vartheta_A\right)
  \label{eq:ac-pr} \\[0.3em]
  Q_R &= \frac{U_S \cdot U_R}{\left|B\right|}
         \cdot \sin\left(\vartheta_B - \delta\right)
       - \frac{\left|A\right| \cdot U_R^2}{\left|B\right|}
         \cdot \sin\left(\vartheta_B - \vartheta_A\right)
  \label{eq:ac-qr}
\end{align}
with $U_S$ and $U_R$ line-to-line magnitudes, so that $P_R$ and $Q_R$ are
three-phase quantities.

Two numbers carry the whole behaviour of the line. The angle $\vartheta_B$
sets where the maximum of $P_R$ falls: differentiating~\eqref{eq:ac-pr} puts
it at $\delta = \vartheta_B$, which is $90^\circ$ only for a lossless line.
The magnitude $\left|A\right|$ scales the second term, a constant offset
independent of $\delta$ that exists only because the line charges itself.

\subsubsection{The short, lossless line}
\label{sec:ac-simplified}

For a short line the shunt admittance vanishes and $\underline{A} =
\underline{D} = 1$, so $\left|A\right| = 1$ and $\vartheta_A = 0$. For a
lossless line $\underline{B} = \underline{Z} = j \cdot X$, so
$\left|B\right| = X$ and $\vartheta_B = 90^\circ$. With
$\cos\left(90^\circ - \delta\right) = \sin\delta$ and $\cos 90^\circ = 0$,
equations~\eqref{eq:ac-pr} and~\eqref{eq:ac-qr} collapse to
\begin{equation}
  P = \frac{U_S \cdot U_R}{X} \cdot \sin\delta
  \qquad
  Q_R = \frac{U_S \cdot U_R \cdot \cos\delta - U_R^2}{X}
  \label{eq:ac-simple}
\end{equation}
with the series reactance
\begin{equation}
  X = \omega \cdot L' \cdot l
  \label{eq:ac-x}
\end{equation}

For small angles this separates further into
$P \approx U_S \cdot U_R \cdot \delta / X$ and
$Q_R \approx U_R \cdot \left(U_S - U_R\right) / X$: active power is carried
by the angle between the terminal voltages, reactive power by the difference
in their magnitude. That decoupling is what makes~\eqref{eq:ac-simple} worth
keeping as a mental model even where it is not accurate enough to compute
with.

The $X$ of~\eqref{eq:ac-simple} is not the line alone but the whole path
between the two internal voltages --- the source impedances and the
transformers at both ends are in series with it --- so the short-circuit
level at the terminals belongs in any statement of transfer capability.

\subsubsection{Distributed representation}
\label{sec:ac-wave}

The $\pi$-section lumps parameters that are in fact distributed along the
line. The distributed description uses two quantities,
\begin{equation}
  \underline{\gamma} = \alpha + j \cdot \beta
    = \sqrt{\underline{Z}' \cdot \underline{Y}'}
  \qquad
  \underline{Z}_w = \sqrt{\frac{\underline{Z}'}{\underline{Y}'}}
  \label{eq:ac-gamma-zw}
\end{equation}
the propagation constant and the wave impedance. Voltage and current are
then defined at every cross-section of the line rather than only at its
terminals, Figure~\ref{fig:ac-terminals}. At a distance $x$ from the
sending end of a line of length $l$ they are
\begin{align}
  \underline{U} &= \underline{U}_R \cdot
      \cosh\left(\underline{\gamma} \cdot \left(l-x\right)\right)
    + \underline{Z}_w \cdot \underline{I}_R \cdot
      \sinh\left(\underline{\gamma} \cdot \left(l-x\right)\right)
  \label{eq:ac-u-x} \\[0.3em]
  \underline{I} &= \frac{\underline{U}_R}{\underline{Z}_w} \cdot
      \sinh\left(\underline{\gamma} \cdot \left(l-x\right)\right)
    + \underline{I}_R \cdot
      \cosh\left(\underline{\gamma} \cdot \left(l-x\right)\right)
  \label{eq:ac-i-x}
\end{align}

\begin{figure}[!ht]
  \centering
  \begin{tikzpicture}[
      x=0.62cm, y=0.62cm, line width=0.8pt,
      every node/.style={inner sep=1.5pt, font=\small},
      rail/.style={atFg, line width=0.8pt},
      term/.style={circle, draw=atFg, fill=atBg, line width=0.8pt,
                   minimum size=4.5pt, inner sep=0pt},
      flow/.style={atFg, -{Triangle[length=2.4mm,width=2.0mm]},
                   line width=0.8pt},
      dim/.style={atFgMuted, dotted, line width=0.9pt},
      ext/.style={atFgMuted, {Stealth[length=1.6mm]}-{Stealth[length=1.6mm]},
                  line width=0.6pt},
    ]
    \def\xa{0}\def\xb{9}\def\xm{3.3}\def\yt{3.6}\def\yb{0}
    % conductors
    \draw[rail] (\xa,\yt) -- (\xb,\yt);
    \draw[rail] (\xa,\yb) -- (\xb,\yb);
    % line current at the cross-section
    \draw[flow] (\xa,\yt) -- (\xm,\yt);
    \node[above=1pt] at (\xm,\yt) {$\underline{I}$};
    % terminals
    \node[term] at (\xa,\yt) {};  \node[term] at (\xb,\yt) {};
    \node[term] at (\xa,\yb) {};  \node[term] at (\xb,\yb) {};
    % terminal currents
    \draw[flow] ({\xa-1.7},\yt) -- ({\xa-0.8},\yt);
    \node[left] at ({\xa-1.9},\yt) {$\underline{I}_S$};
    \draw[flow] ({\xb+0.8},\yt) -- ({\xb+1.7},\yt);
    \node[right] at ({\xb+1.9},\yt) {$\underline{I}_R$};
    % voltages, upper conductor to lower conductor
    \draw[flow] (\xa,{\yt-0.4}) -- (\xa,{\yb+0.4});
    \node[right=2pt] at (\xa,{(\yt+\yb)/2}) {$\underline{U}_S$};
    \draw[flow] (\xm,{\yt-0.4}) -- (\xm,{\yb+0.4});
    \node[right=2pt] at (\xm,{(\yt+\yb)/2}) {$\underline{U}$};
    \draw[flow] (\xb,{\yt-0.4}) -- (\xb,{\yb+0.4});
    \node[right=2pt] at (\xb,{(\yt+\yb)/2}) {$\underline{U}_R$};
    % coordinate x and total length l
    \draw[dim] (\xa,{\yb-0.3}) -- (\xa,{\yb-2.6});
    \draw[dim] (\xm,{\yb-0.3}) -- (\xm,{\yb-1.6});
    \draw[dim] (\xb,{\yb-0.3}) -- (\xb,{\yb-2.6});
    \draw[ext] (\xa,{\yb-1.2}) -- (\xm,{\yb-1.2});
    \node[below=1pt, atFgMuted] at ({(\xa+\xm)/2},{\yb-1.2}) {$x$};
    \draw[ext] (\xa,{\yb-2.2}) -- (\xb,{\yb-2.2});
    \node[below=1pt, atFgMuted] at ({(\xa+\xb)/2},{\yb-2.2}) {$l$};
  \end{tikzpicture}
  \caption{Terminal quantities of the two-port and the local voltage and
           current at a distance $x$ along the line}
  \label{fig:ac-terminals}
\end{figure}

Evaluating~\eqref{eq:ac-u-x} and~\eqref{eq:ac-i-x} at $x = 0$ and comparing
with~\eqref{eq:ac-abcd} identifies the exact two-port parameters,
\begin{equation}
  \underline{A} = \underline{D} = \cosh\left(\underline{\gamma} \cdot l\right),
  \qquad
  \underline{B} = \underline{Z}_w \cdot
      \sinh\left(\underline{\gamma} \cdot l\right),
  \qquad
  \underline{C} = \frac{\sinh\left(\underline{\gamma} \cdot l\right)}
                       {\underline{Z}_w}
  \label{eq:ac-abcd-exact}
\end{equation}
which replace the nominal-$\pi$ entries of~\eqref{eq:ac-abcd-pi}. The
relations~\eqref{eq:ac-pr} and~\eqref{eq:ac-qr} are unchanged; only the way
$\underline{A}$ and $\underline{B}$ are computed differs.

The exact result can itself be written as a $\pi$. Matching the branches of a
$\pi$-section to~\eqref{eq:ac-abcd-exact} entry by entry gives
\begin{equation}
  \underline{Z}_{eq} = \underline{Z}_w \cdot
      \sinh\left(\underline{\gamma} \cdot l\right),
  \qquad
  \frac{\underline{Y}_{eq}}{2} = \frac{1}{\underline{Z}_w} \cdot
      \tanh\left(\frac{\underline{\gamma} \cdot l}{2}\right)
  \label{eq:ac-equivalent-pi}
\end{equation}
the \emph{equivalent} $\pi$, often called the long-line correction. It is not
an approximation: one $\pi$ with these two branches reproduces the
distributed line exactly at its terminals at that frequency, because
$1 + \sinh\left(\underline{\gamma} \cdot l\right) \cdot
\tanh\left(\underline{\gamma} \cdot l / 2\right) =
\cosh\left(\underline{\gamma} \cdot l\right)$. The nominal $\pi$
of~\eqref{eq:ac-abcd-pi} is its small-argument limit, with
$\sinh u \to u$ and $\tanh u \to u$. This is what a load-flow program
stores for a long line, and it is why a line keeps a $\pi$-shaped model
however long it is. Section~\ref{sec:ref-cascade} puts numbers on both.

What decides whether the lumped form will do is the electrical length
$\beta \cdot l$, not the length in km. Since $\beta = \omega \cdot
\sqrt{L' \cdot C'}$ for a low-loss line, a cable reaches a given electrical
length in a fraction of the distance an overhead line needs, and the lumped
approximation fails correspondingly earlier.

\FloatBarrier
